\Needspace{16\baselineskip}
\Article{356}
\textsc{Theorem.} \emph{If $P$ is a prime function of $x$ of degree $m$, and $X$ is a function of $x,x^p,x^{pp},x^{p^3}\dotstwo x^{p^{m-1}}$ in which all these quantities occur in the same way, that is, which remains unchanged however they are permuted among themselves, then dividing the function $X$ by $P$ gives a remainder which is a number.}

\emph{Proof.} Let the remainder be
\[
Ax^{m-1}+Bx^{m-2}+\dotstwo+N\equiv\xi
\]
All coefficients $A,B,C\ldots$ up to, but excluding, $N$ will be $\equiv0$. This is proved as follows. Since $X-\xi$ is divisible by $P$, $X^p-\xi^p$ can also be divided by $P$.
\SourceLeaf{238}
But it is easy to see that $X^p$ is what $X$ becomes when $x$ is replaced by $x^p$, $x^p$ by $x^{pp}$, and so on, .\ and $x^{p^{m-1}}$ by $x^{p^m}$, or, equivalently, $x$. Thus it is clear that $X^p\equiv X\;(\mathrm{mod.}\,P)$; hence, since $X^p\equiv\xi^p$ and $X\equiv\xi(\mathrm{mod.}\,P)$, we also have $\xi^p\equiv\xi(\mathrm{mod.}\,P$, or
\[
\xi^p-\xi\equiv0(\mathrm{mod.}\,P)
\]
But $\xi^p-\xi$ is congruent, modulo $p$, to the product of $\xi,\xi+1,\xi+2,\dotstwo$ through $\xi+p-1$. These factors are all coprime to $P$ unless $\xi$ is simply a number. Hence $\xi^p-\xi$ cannot be divisible by $P$ in any other way. Q. E. D.

Functions of this kind are the sum of all the quantities, the sum of their squares, cubes, and so on, and the sums of their products taken two at a time, three at a time, and so on. We shall determine what that number is in the following article:

\clearpage
\Article{357}
\textsc{Theorem.} \emph{Let the prime function of the preceding article be}
\[
P\equiv x^m-Ax^{m-1}+Bx^{m-2}-Cx^{m-3}+\text{etc.}
\]
\emph{Then the remainder when the sum of the quantities $x,x^p$, and so on, $x^{p^{m-1}}$ is divided by $P$ is $\equiv A$; for the sum of their products taken two at a time it is $\equiv B$; for the sum of their products taken three at a time it is $\equiv C$, and so on.}

\emph{Proof.} Let those functions be $X,Y,Z$, and so on, and let their remainders be the numbers $A',B',C'$, and so on, in their respective order. It is easy to see that $x,x^p,x^{pp}$, and so on, are roots of the equation
\[
z^m-Xz^{m-1}+Yz^{m-2}-Zz^{m-3}+\text{etc.}=0
\]
Thus, on putting $z=x$, we have
\[
x^m-Xx^{m-1}+Yx^{m-2}-Zx^{m-3}+\text{etc.}=0
\]
But the functions $X-A'$, $Y-B'$, $Z-C'$, and so on, are divisible by $P$, and hence so is the function
\[
x^m-A'x^{m-1}+B'x^{m-2}-C'x^{m-3}+\text{etc.}
\]
This cannot happen unless $A'\equiv A$, $B'\equiv B$, $C'\equiv C$, and so on. Q. E. D.

It is also known that any other function $X$ of $x,x^p$, and so on, [in which all these quantities occur in the same way,] can always be derived from these. Thus
\SourceLeaf{239}
\[
x^2+x^{2p}+x^{3pp}+\text{etc.}\equiv AA-2B(\mathrm{mod.}\,P)\text{ etc. etc.}
\]
Example. Let $p=5$ and $P\equiv x^2+2x+3$. Then the function $x+x^5$, divided by $P$, gives $\equiv-2$, and $x^6\equiv3$, and so on, and so on.
